\documentclass[12pt,a4paper]{article}
\usepackage[utf8]{inputenc}
\usepackage[T1]{fontenc}
\usepackage[greek,english]{babel}
\usepackage{geometry}
\usepackage{booktabs}
\usepackage{amsmath}
\usepackage{hyperref}
\usepackage{setspace}
\usepackage{array}

\geometry{margin=2.5cm}
\onehalfspacing

\hypersetup{colorlinks=true, linkcolor=blue, urlcolor=blue, citecolor=blue}

\title{\selectlanguage{greek}Στατιστική Ανάλυση του Φαιστίου Δίσκου:\\
Υπολογιστική Μεθοδολογία Αξιολόγησης Φωνητικών Κλειδιών}

\author{\selectlanguage{greek}Μανόλης Χαβαδάκης\\
\small \selectlanguage{greek}Ανεξάρτητος Ερευνητής\\
\small \href{mailto:aikopappas@gmail.com}{aikopappas@gmail.com}}

\date{\selectlanguage{greek}Ιούνιος 2026 \quad Έκδοση 3.3}

\begin{document}
\selectlanguage{greek}
\maketitle

\begin{abstract}
Ο Φαίστιος Δίσκος ($\sim$1700 π.Χ.) παραμένει ένα από τα πλέον αμφιλεγόμενα αναλυθέντα αντικείμενα της αρχαιολογίας. Παρουσιάζουμε ένα τυφλό υπολογιστικό πλαίσιο αξιολόγησης ανταγωνιστικών φωνητικών κλειδιών, εφαρμόζοντας Bonferroni-διορθωμένη Monte Carlo προσομοίωση σε 10 υποψήφια κλειδιά έναντι τριών σωμάτων αναφοράς: Λουβικό Ιερογλυφικό λεξιλόγιο, πίνακες συχνοτήτων Γραμμικής~Α, και το Αιγυπτιακό σώμα AED-TEI (675.773 tokens από 13.950 κείμενα). Τρία κλειδι-ανεξάρτητα ευρήματα τεκμηριώνονται ανεξαρτήτως φωνητικών υποθέσεων: (1)~το διγράμματο [Σημείο\#36$\to$Σημείο\#11] εμφανίζει $Z=10$ πλεονάζουσα γειτνίαση (παρατ./αναμ.$=7{,}69\times$, $p\approx 0$)· (2)~ο έλεγχος τομέα σώματος επιβεβαιώνει θεολογική ταξινόμηση ($Z=27{,}16$ έναντι διοικητικού $Z=-0{,}40$)· (3)~το Σημείο~\#45 (ηλιακός ρόδακας) εμφανίζεται αποκλειστικά στα κέντρα A31 και B30. Το Λουβικό Ιερογλυφικό κλειδί (\textsc{G\_Luwian}) επιτυγχάνει την υψηλότερη Bonferroni-σημαντική βαθμολογία ($S=523$, $p<0{,}0001$), παράγοντας ηλιακο-υδάτινη κοσμολογική ανάγνωση δομικά παράλληλη με το Αιγυπτιακό \textit{Amduat}. Αρνητικό τεστ ελέγχου σε συνθετικό δίσκο ($Z=1{,}99$, μη σημαντικό) αποκαλύπτει ότι η βαθμολογία token-επιπέδου είναι $\sim$94\% συχνοτητο-εξαρτώμενη. Μόνο οι τρεις κλειδι-ανεξάρτητοι πυλώνες παρουσιάζονται ως πρωτεύοντες ισχυρισμοί. Ο κώδικας δημοσιεύεται ανοιχτά.

\medskip
\noindent\textbf{Λέξεις-κλειδιά:} Φαίστιος Δίσκος, αδιάγνωστα συστήματα γραφής, υπολογιστική γλωσσολογία, Λουβικά ιερογλυφικά, Monte Carlo, Bonferroni, Αιγαίο Εποχής Χαλκού
\end{abstract}

\tableofcontents
\newpage

%------------------------------------------------------------
\section{Εισαγωγή}
%------------------------------------------------------------

Ο Φαίστιος Δίσκος, που ανακαλύφθηκε το 1908 στο Μινωικό ανάκτορο της Φαιστού (Κρήτη) και χρονολογείται περίπου στο 1700~π.Χ., φέρει 241 αποτυπωμένα σημεία από 45 διακριτά σύμβολα σε αμφίπλευρη σπείρα (61 ομάδες λέξεων). Παραμένει μοναδικό εύρημα χωρίς δεύτερο παράδειγμα.

Οι προηγούμενες απόπειρες αποκρυπτογράφησης αριθμούν εκατοντάδες και μοιράζονται κοινή μεθοδολογική αδυναμία: το προτεινόμενο κλειδί κατασκευάζεται ώστε να παράγει εύλογες αναγνώσεις, δημιουργώντας ανελέγκτη κυκλικότητα.

Η παρούσα μελέτη \textbf{δεν προτείνει αποκρυπτογράφηση}. Προτείνει \textbf{στατιστική μεθοδολογία} κατάταξης ανταγωνιστικών φωνητικών κλειδιών έναντι αντικειμενικών σωμάτων αναφοράς, με ρητή διόρθωση για πολλαπλές συγκρίσεις.

%------------------------------------------------------------
\section{Προηγούμενη Έρευνα}
%------------------------------------------------------------

\textbf{Γλωσσολογικές προτάσεις} (Owens 1996· Achterberg κ.ά. 2004· Weingarten 2016): αναθέτουν φωνητικές αξίες χωρίς στατιστική επικύρωση. \textbf{Δομικές αναλύσεις} (Sproat 2010· Rao κ.ά. 2009): επιβεβαιώνουν σύστημα γραφής χωρίς να προσδιορίζουν γλώσσα. \textbf{Υπολογιστικές προσεγγίσεις}: καμία προηγούμενη μελέτη δεν εφαρμόζει Bonferroni-διορθωμένη Monte Carlo σε πολλαπλά κλειδιά έναντι ανεξάρτητων σωμάτων αναφοράς.

%------------------------------------------------------------
\section{Δεδομένα}
%------------------------------------------------------------

\subsection{Ο Φαίστιος Δίσκος}

Κλειδι-ανεξάρτητες παρατηρήσεις:
\begin{itemize}
  \item Σημείο~\#2 είναι αρχή-λέξης σε 48\% των λέξεων
  \item Σημείο~\#11 είναι τέλος-λέξης σε 30\% των λέξεων
  \item Shannon εντροπία $H = 3{,}045$ bits (συλλαβικό: 2,8--3,5~bits)
  \item Zipf $R^2 = 0{,}673$ — τελετουργική/φορμουλαϊκή γλώσσα
\end{itemize}

\subsection{Σώματα Αναφοράς}

\paragraph{AED-TEI Αιγυπτιακό Σώμα.}
675.773 tokens, 13.950 κείμενα. Τελετουργικό υποσώμα: 95.162 tokens, 1.370 κείμενα (Κείμενα Πυραμίδων, Βίβλος Νεκρών, Κείμενα Σαρκοφάγων). Άδεια: CC-BY-SA~4.0.

\paragraph{Λουβικό Ιερογλυφικό Λεξιλόγιο (G\_LUWIAN).}
19 καταχωρήσεις με τεκμηριωμένες φωνητικές αξίες (Masson 1961· Hawkins 2000· Melchert 2003). Ανατολιακές επιγραφές $\sim$2000--1200~π.Χ.

\paragraph{Πίνακας Συχνοτήτων Γραμμικής Α (B\_FREQ).}
30 αντιστοιχίες σημείου-συλλαβής. Η Γραμμική Α παραμένει αδιάγνωστη.

%------------------------------------------------------------
\section{Μεθοδολογία}
%------------------------------------------------------------

\subsection{Τυφλός Δοκιμαστικός Πλέγμας}
Δέκα κλειδιά (Α--J) κατασκευάστηκαν ανεξάρτητα και αξιολογήθηκαν ταυτόχρονα χωρίς εκ των υστέρων τροποποίηση.

\subsection{Συνάρτηση Βαθμολόγησης}
\begin{equation}
\text{score}(W, K) = \sum_{i=1}^{n} \max\_\text{length\_match}\!\left(K(s_i),\, \mathcal{V}\right)
\end{equation}
Συνολικό $S(K) = \sum_W \text{score}(W, K)$.

\subsection{Κατανομή Null Monte Carlo}
Κλειδί~J: 10.000 τυχαίες αντιστοιχίσεις (seed~=~42):
\[
\mu_\text{null} = 151{,}9, \quad \sigma_\text{null} = 77{,}0
\]
Κατώφλι Bonferroni ($p < 0{,}005$, διορθ. για 10~κλειδιά): $S > 379$.
Κατώφλι δημοσίευσης ($p < 0{,}0001$): $S > 521$.

\subsection{Έλεγχος Τομέα Σώματος}
Λεξιλόγιο G\_LUWIAN βαθμολογήθηκε χωριστά έναντι θεολογικού (Κείμενα Πυραμίδων κ.λπ.) και διοικητικού υποσώματος AED-TEI.

\subsection{Αρνητικός Έλεγχος}
1.000 συνθετικοί δίσκοι: πανομοιότυπη κατανομή συχνοτήτων, τυχαιοποιημένη γειτνίαση.

%------------------------------------------------------------
\section{Αποτελέσματα}
%------------------------------------------------------------

\subsection{Κατάταξη Κλειδιών}

\begin{table}[h]
\centering
\caption{Bonferroni-διορθωμένη κατάταξη κλειδιών. Κατώφλι: $S > 379$ ($p < 0{,}005$).}
\begin{tabular}{llrrrrl}
\toprule
Θέση & Κλειδί & Βαθμ. & $Z$ & $p$-τιμή & Bonf. & Refrain \\
\midrule
1 & \textsc{G\_Luwian} & \textbf{523} & 4,82 & \textbf{$<$0,0001} & $\checkmark\checkmark\checkmark$ & \textit{za-wa-tar} \\
2 & \textsc{E1\_Egypt}  & 491 & 4,40 & 0,0001  & $\checkmark\checkmark$ & \textit{n-m-r} \\
3 & \textsc{B\_Freq}    & 430 & 3,61 & 0,0009  & $\checkmark\checkmark$ & \textit{a-sa-ra} \\
4 & \textsc{I\_Morpho}  & 426 & 3,56 & 0,0009  & $\checkmark\checkmark$ & \textit{a-ku-te} \\
5 & \textsc{E2\_Wsir}   & 261 & 1,42 & 0,09    & ---  & \textit{A-sa-r} \\
6 & \textsc{A\_Evans}   & 188 & 0,47 & 0,30    & ---  & --- \\
7 & \textsc{F\_Cypriot} & 156 & 0,05 & 0,45    & ---  & \textit{a-ku-se} \\
8 & \textsc{H\_Abjad}   & 0   & $-1{,}97$ & 1,00 & --- & \textsc{αποκλείεται} \\
\bottomrule
\end{tabular}
\end{table}

\subsection{Τρεις Κλειδι-Ανεξάρτητοι Πυλώνες}

\paragraph{Πυλώνας~1 --- Διγράμματο [Σημείο\#36$\to$Σημείο\#11] ($p \approx 0$).}
Παρατηρούμενα: 17. Αναμενόμενα: 2,2. Λόγος: $7{,}69\times$. $Z = 10{,}00$. Υπό \textsc{G\_Luwian}: \textit{wa-tar} (ΠΙΕ~*\textit{wódr̥}). Δεν εξηγείται από οριακές συχνότητες.

\paragraph{Πυλώνας~2 --- Έλεγχος τομέα σώματος ($Z=27{,}16$).}
Θεολογικό υποσώμα: $Z = 27{,}16$. Διοικητικό: $Z = -0{,}40$. Διαφορά $27{,}5\sigma$ επιβεβαιώνει τελετουργική κατηγοριοποίηση ανεξαρτήτως κλειδιού.

\paragraph{Πυλώνας~3 --- Σημείο~\#45 μόνο στα κέντρα.}
Εμφανίζεται αποκλειστικά σε A31 και B30. Πιθανότητα τυχαίας εμφάνισης: $p \approx 0{,}0018$.

\subsection{Χιασμός Κέντρων}
\begin{align*}
\text{A31} &= [45,\; 2,\; 36,\; 11,\; 22] \\
\text{B30} &= [45,\; 36,\; 11,\; 2,\; 6]
\end{align*}
Κοινά σημεία $\{2, 11, 36, 45\}$ σε αντεστραμμένη σειρά δεικτών. Κλειδι-ανεξάρτητη γεωμετρική ιδιότητα.

\subsection{Ευαισθησία και Διασταύρωση}
105/105 διαταραχές άνω Bonferroni. Spearman $\rho = 0{,}779$. Μεταφορά Α$\to$Β: 64,6\%· Β$\to$Α: 91,6\%.

%------------------------------------------------------------
\section{Αρνητικός Έλεγχος και Αυτοκριτική}
%------------------------------------------------------------

\subsection{Συχνοτητο-Εξαρτώμενη Βαθμολογία Token}
Μέσο $Z$ σε συνθετικούς δίσκους: $1{,}99$ (μη σημαντικό). Περίπου 94\% της βαθμολογίας token-επιπέδου εξηγείται από οριακές συχνότητες. \textbf{Η βαθμολογία $Z=8{,}58$ αποσύρεται ως κύριο επιχείρημα.}

\subsection{Κυκλικότητα Σχεδιασμού}
\textsc{G\_Luwian} κατασκευάστηκε με γνώση στατιστικών δίσκου. Κύρια αδιόρθωτη αδυναμία. Θεραπεία: τυφλή αναπαραγωγή από Λουβιολόγο.

%------------------------------------------------------------
\section{Συζήτηση}
%------------------------------------------------------------

\subsection{Κύρια Ερμηνεία}
Υπό \textsc{G\_Luwian}: refrain~$=$ \textit{za-wa-tar} («αυτό το νερό», ΠΙΕ~*\textit{wódr̥})· κέντρο A31~$=$ \textit{ti-wa-za-wa-tar-ha} («TIWAT! αυτό το νερό --- ναι!»)· κέντρο B30~$=$ \textit{ti-wa-wa-tar-za-an} («TIWAT! νερό-κριτής --- εδώ!»). Ηλιακο-υδάτινος κοσμολογικός ύμνος: ο Tiwat κατεβαίνει στα πρωτόγονα νερά (Πλευρά~Α) και ανεβαίνει αναγεννημένος (Πλευρά~Β).

\subsection{Αιγυπτιακό Δομικό Παράλληλο}
Το \textit{Amduat} περιγράφει τη νυχτερινή σπειροειδή κάθοδο του Ra στο Nun, ένωση με τον Osiris τα μεσάνυχτα, και ηλιακή αναγέννηση. Δομική αντιστοίχιση: Ra~$\leftrightarrow$~Tiwat· Nun/Osiris~$\leftrightarrow$~\textit{watar}· μεσάνυχτο~$\leftrightarrow$~κέντρα A31/B30· κάθοδος~$\leftrightarrow$~Πλευρά~Α· άνοδος~$\leftrightarrow$~Πλευρά~Β.

%------------------------------------------------------------
\section{Περιορισμοί}
%------------------------------------------------------------
\begin{enumerate}
  \item \textbf{Κυκλικότητα:} G\_Luwian κατασκευάστηκε γνωρίζοντας τα στατιστικά. Διερευνητικό.
  \item \textbf{Hapax legomenon:} κανένα δεύτερο Φαίστιο κείμενο.
  \item \textbf{Βαθμολογία token:} $\sim$94\% συχνοτητο-εξαρτώμενη. Αποσύρεται.
  \item \textbf{Λεξιλόγιο:} G\_Luwian καλύπτει μόνο 19 καταχωρήσεις.
  \item \textbf{Γραμμική Α:} B\_Freq παρεκτείνει από αδιάγνωστο σύστημα.
\end{enumerate}

%------------------------------------------------------------
\section{Συμπεράσματα}
%------------------------------------------------------------

Αποδείξαμε: (1)~μη τυχαίο διγράμματο [Σημ.\#36$\to$Σημ.\#11] ($Z=10$)· (2)~αναγνωρίσιμη τελετουργική γλωσσική ποικιλία ($Z=27{,}16$ έναντι $Z=-0{,}40$)· (3)~γεωμετρικό χιασμό και Σημείο~\#45 αποκλειστικά στα κέντρα· (4)~\textsc{G\_Luwian} με $p<0{,}0001$· (5)~συνεπή ηλιακο-υδάτινη ανάγνωση με παράλληλα \textit{Amduat}· (6)~βαθμολογία token $\sim$94\% συχνοτητο-εξαρτώμενη.

\textbf{Η ανεξάρτητη αναπαραγωγή από Λουβιολόγο παραμένει το κρίσιμο επόμενο βήμα.}

%------------------------------------------------------------
\begin{thebibliography}{99}
\selectlanguage{english}
%------------------------------------------------------------
\bibitem{achterberg2004} Achterberg W., Best J., Enzler K., Rietveld L. \& Woudhuizen F. (2004). \textit{The Phaistos Disc: A Luwian Letter to Nestor}. Dutch Monographs on Ancient History and Archaeology.
\bibitem{assmann2001} Assmann J. (2001). \textit{The Search for God in Ancient Egypt}. Cornell University Press.
\bibitem{faulkner1969} Faulkner R.O. (1969). \textit{The Ancient Egyptian Pyramid Texts}. Oxford University Press.
\bibitem{hawkins2000} Hawkins J.D. (2000). \textit{Corpus of Hieroglyphic Luwian Inscriptions}. De Gruyter.
\bibitem{hornung1999} Hornung E. (1999). \textit{The Ancient Egyptian Books of the Afterlife}. Cornell University Press.
\bibitem{masson1961} Masson E. (1961). \textit{Recherches sur les plus anciens emprunts sémitiques en grec}. Paris.
\bibitem{melchert2003} Melchert H.C. (2003). \textit{The Luwians}. Brill.
\bibitem{owens1996} Owens G. (1996). The Phaistos Disc: A New Approach. \textit{Cretan Studies}, 5, 1--24.
\bibitem{rao2009} Rao R.P.N. et al. (2009). Entropic Evidence for Linguistic Structure in the Indus Script. \textit{Science}, 324, 1165.
\bibitem{schweitzer2011} Schweitzer S.D. (2011). AED-TEI. \url{https://github.com/simondschweitzer/aed-tei}
\bibitem{sproat2010} Sproat R. (2010). Ancient Symbols, Computational Linguistics... \textit{Computational Linguistics}, 36(3), 585--594.
\bibitem{weingarten2016} Weingarten J. (2016). The Phaistos Disc: Pedigree of a Forgery. \textit{Journal of Prehistoric Religion}, 25.
\end{thebibliography}

\end{document}
